\chapter{Pointing reduces to one scalar}
\label{chap:pointing}

With yaw free, whether the hinged camera can see a direction depends on a single
number: the elevation of that direction above the rotor plane. Throughout this
chapter $b_3$ is a unit thrust axis that is not horizontal
($\inner{b_3}{\up}\neq0$) and $u$ is a unit line of sight.

\begin{definition}[Rotor-plane elevation]
\label{def:elevation}
\lean{VisionTracking.elevation}
\leanok
$e(b_3,u)=\arcsin\inner{b_3}{u}\in[-\pi/2,\pi/2]$.
\end{definition}

\begin{lemma}[Orthonormal-pair bound]
\label{lem:orthonormal-bound}
\lean{VisionTracking.inner_sq_add_inner_sq_le, VisionTracking.inner_camera_le_cos,
  VisionTracking.abs_sub_le_angle_of_orthonormal, VisionTracking.angle_eq_abs_sub_of_aligned}
\leanok
Let $b_1\perp b_3$ be unit vectors and $\eta\in[-\pi/2,\pi/2]$. Then
$\ang(\cos\eta\,b_1+\sin\eta\,b_3,\ u)\ge|\eta-\arcsin\inner{b_3}{u}|$, with
equality when $\inner{b_1}{u}=\sqrt{1-\inner{b_3}{u}^2}$.
\end{lemma}

\begin{proof}
\leanok
Bessel's inequality gives $\inner{b_1}{u}^2+\inner{b_3}{u}^2\le1$, so
$\inner{b_1}{u}\le\cos e$. Since $\cos\eta\ge0$,
$\inner{c}{u}\le\cos\eta\cos e+\sin\eta\sin e=\cos(\eta-e)$. Apply the decreasing
function $\arccos$ and use $|\eta-e|\le\pi$.
\end{proof}

\begin{lemma}[Body forward axis]
\label{lem:body-forward}
\lean{VisionTracking.norm_bodyForward, VisionTracking.inner_bodyForward_self,
  VisionTracking.heading_proj_ne_zero}
\leanok
\uses{def:cam-axis}
$b_1=\mathrm{bodyForward}(b_3,\psi)$ is a unit vector orthogonal to $b_3$.
\end{lemma}

\begin{proof}
\leanok
The heading is horizontal and of unit length. Its projection onto $b_3^\perp$ could
only vanish if the heading were a multiple of $b_3$, and taking the inner product
with $\up$ rules that out.
\end{proof}

\begin{lemma}[Yaw covers the rotor plane]
\label{lem:yaw-covers}
\lean{VisionTracking.exists_yaw_bodyForward_eq}
\leanok
\uses{lem:body-forward}
Every unit vector $w\perp b_3$ equals $\mathrm{bodyForward}(b_3,\psi)$ for some
yaw $\psi$.
\end{lemma}

\begin{proof}
\leanok
Slide $w$ along $b_3$ to the horizontal vector
$h=w-\frac{\inner{w}{\up}}{\inner{b_3}{\up}}b_3$, which is nonzero. Take $\psi$ to
be the argument of $(h_0,h_1)$, so that the heading at $\psi$ is $h/\norm{h}$. Its
projection onto $b_3^\perp$ is $w/\norm{h}$, which normalises to $w$.
\end{proof}

\begin{theorem}[Pointing reduction]
\label{thm:pointing}
\lean{VisionTracking.abs_sub_elevation_le_angle, VisionTracking.exists_yaw_angle_eq,
  VisionTracking.exists_visible_iff}
\leanok
\uses{def:elevation, lem:orthonormal-bound, lem:yaw-covers}
For every yaw $\psi$ and every $\eta\in[-\pi/2,\pi/2]$,
$\ang(c(b_3,\psi,\eta),u)\ge|\eta-e(b_3,u)|$, and some yaw attains equality.
Consequently, for a hinge range
$-\pi/2\le\eta_{\mathrm{lo}}\le\eta_{\mathrm{hi}}\le\pi/2$ and $\alpha\ge0$,
\[
  \exists\,\psi,\ \eta\in[\eta_{\mathrm{lo}},\eta_{\mathrm{hi}}]:\
  \ang(c(b_3,\psi,\eta),u)\le\alpha
  \iff
  \eta_{\mathrm{lo}}-\alpha\le e(b_3,u)\le\eta_{\mathrm{hi}}+\alpha .
\]
\end{theorem}

\begin{proof}
\leanok
The lower bound is Lemma~\ref{lem:orthonormal-bound} applied with
Lemma~\ref{lem:body-forward}. For attainment, align $b_1$ with the rotor-plane
part of $u$ using Lemma~\ref{lem:yaw-covers}; when $u=\pm b_3$ every yaw is
optimal. For the equivalence, take $\eta$ to be $e$ clamped to the hinge range.
\end{proof}

\section{Overhead geometry}

When the drone hovers directly above the target, the line of sight is $-\up$.

\begin{lemma}[Overhead elevation]
\label{lem:overhead-elevation}
\lean{VisionTracking.elevation_neg_up}
\leanok
\uses{def:elevation}
For a unit $b_3$, $e(b_3,-\up)=\ang(b_3,\up)-\pi/2$: the elevation of ``straight
down'' is the tilt minus a right angle.
\end{lemma}

\begin{proof}
\leanok
$\inner{b_3}{-\up}=-\cos t$ with $t=\ang(b_3,\up)\in[0,\pi]$, and
$\arcsin(-\cos t)=-\arcsin\sin(\pi/2-t)=t-\pi/2$.
\end{proof}

\begin{theorem}[Overhead visibility across a tilt band]
\label{thm:overhead-band}
\lean{VisionTracking.overhead_visible_iff, VisionTracking.overhead_band_iff,
  VisionTracking.TiltBand, VisionTracking.tiltVec}
\leanok
\uses{thm:pointing, lem:overhead-elevation}
Let $0\le t_{\mathrm{lo}}\le t_{\mathrm{hi}}<\pi/2$. The camera, with yaw free and
hinge in $[\eta_{\mathrm{lo}},\eta_{\mathrm{hi}}]$, can look straight down for
\emph{every} unit thrust axis with tilt in $[t_{\mathrm{lo}},t_{\mathrm{hi}}]$ iff
\[
  \eta_{\mathrm{lo}}\le t_{\mathrm{lo}}-\tfrac{\pi}{2}+\alpha
  \quad\text{and}\quad
  \eta_{\mathrm{hi}}\ge t_{\mathrm{hi}}-\tfrac{\pi}{2}-\alpha .
\]
\end{theorem}

\begin{proof}
\leanok
Sufficiency follows from Theorem~\ref{thm:pointing} and
Lemma~\ref{lem:overhead-elevation}. For necessity, test the two extreme axes
$\cos t\,\up+\sin t\,e_1$ with $t=t_{\mathrm{lo}}$ and $t=t_{\mathrm{hi}}$.
\end{proof}

\begin{theorem}[Fixed camera from overhead]
\label{thm:fixed-overhead}
\lean{VisionTracking.fixed_overhead_iff}
\leanok
\uses{thm:overhead-band}
Some fixed mount angle $\eta_0\in[-\pi/2,\pi/2]$ sees straight down at every tilt in
$[t_{\mathrm{lo}},t_{\mathrm{hi}}]$ iff $t_{\mathrm{hi}}-t_{\mathrm{lo}}\le2\alpha$.
The mount $\eta_0=(t_{\mathrm{lo}}+t_{\mathrm{hi}})/2-\pi/2$ then works. In
particular, a camera fixed straight down ($\eta_0=-\pi/2$) needs
$t_{\mathrm{hi}}\le\alpha$.
\end{theorem}

\begin{proof}
\leanok
Apply Theorem~\ref{thm:overhead-band} with
$\eta_{\mathrm{lo}}=\eta_{\mathrm{hi}}=\eta_0$.
\end{proof}

\begin{remark}
In words: a hinge buys exactly the excess of the tilt band over $2\alpha$. If
copying the target forces tilts from $10^\circ$ to $75^\circ$, a $90^\circ$
camera ($\alpha=45^\circ$) works with no hinge at all, mounted $47.5^\circ$ below
the body's forward axis. A $60^\circ$ camera ($\alpha=30^\circ$) needs $5^\circ$ of
hinge travel. The open problem is the tilt band itself, which depends on drag at
the target's speeds and on its manoeuvres (Definition~\ref{def:thrust-demand}).
\end{remark}
